、
4
4\documentclass{Preport}

\usepackage{hyperref}
\hypersetup{
    colorlinks=true,
}
\usepackage{amsmath}\usepackage{float}
\usepackage{circuitikz}
% 双臂电桥（开尔文电桥）原理电路图
% 拓扑：R1(Px1→B)、R2(B→Ps2) 与主回路 Rx-r-Rs 并联；
%       R3(Px2→D)、R4(D→Ps1) 与粗导线 r 并联；检流计 G 接在 B、D 之间。
%       平衡时由 R1 I1 = Rx I3 + R3 I2、R2 I1 = Rs I3 + R4 I2 得 Rx = (R1/R2) Rs。
\newcommand{\ArmBridge}{%
\begin{circuitikz}[european,line width=0.7pt]

  % ---------- 顶部主回路：四端电阻 Rx —— 粗导线 r —— 比较臂 Rs ----------
  \draw (1.5,4.8) rectangle (5.1,5.2);
  \node[font=\small] at (3.3,5.0) {$R_x$};
  \draw (6.4,4.8) rectangle (10.0,5.2);
  \node[font=\small] at (8.2,5.0) {$R_s$};
  \draw[line width=2pt] (5.1,5.0) -- (6.4,5.0);
  \node[below,font=\footnotesize] at (5.75,4.94) {$r$};

  % ---------- 比率臂：R1/R2、R3/R4 各留出间隙 ----------
  \draw (2.1,4.8) -- (2.1,4.4);
  \draw (2.1,4.4) to[R] (2.1,3.4);
  \draw (2.1,3.4) -- (2.1,2.8);
  \draw (2.1,2.8) to[R] (2.1,1.8);
  \draw (4.5,4.8) -- (4.5,4.4);
  \draw (4.5,4.4) to[R] (4.5,3.4);
  \draw (4.5,3.4) -- (4.5,2.8);
  \draw (4.5,2.8) to[R] (4.5,1.8);

  % ---------- 检流计 ----------
  \node[circle,draw,inner sep=1pt,font=\small] (G) at (3.3,3.1) {$G$};
  \draw (2.1,3.1) -- (G.west);
  \draw (G.east) -- (4.5,3.1);
  \fill (2.1,3.1) circle (1pt);
  \fill (4.5,3.1) circle (1pt);

  % ---------- 电位端回路 ----------
  \draw (2.1,1.8) -- (2.1,0.6) -- (9.4,0.6) -- (9.4,4.8);   % R2 → Ps2
  \draw (4.5,1.8) -- (4.5,1.2) -- (7.0,1.2) -- (7.0,4.8);   % R4 → Ps1

  % ---------- 直流电源 ----------
  \draw (1.5,5.0) -- (-0.4,5.0) -- (-0.4,-1.3) -- (4.0,-1.3);
  \draw (4.0,-1.52) -- (4.0,-1.08);
  \draw (4.45,-1.42) -- (4.45,-1.18);
  \draw (4.45,-1.3) -- (10.8,-1.3) -- (10.8,5.0) -- (10.0,5.0);
  \node[below,font=\small] at (4.22,-1.72) {$\varepsilon$};

  % ---------- 端子与元件标注 ----------
  \node[above left,font=\footnotesize]  at (1.45,5.06)  {$C_{x1}$};
  \node[above,font=\footnotesize]       at (4.95,5.22)  {$C_{x2}$};
  \node[above,font=\footnotesize]       at (6.55,5.22)  {$C_{s1}$};
  \node[above right,font=\footnotesize] at (10.05,5.06) {$C_{s2}$};
  \node[right,font=\footnotesize]       at (2.16,4.6)   {$P_{x1}$};
  \node[right,font=\footnotesize]       at (4.56,4.6)   {$P_{x2}$};
  \node[right,font=\footnotesize]       at (7.06,4.6)   {$P_{s1}$};
  \node[right,font=\footnotesize]       at (9.46,4.6)   {$P_{s2}$};
  \node[left,font=\footnotesize]        at (1.86,3.9)   {$R_1$};
  \node[left,font=\footnotesize]        at (1.86,2.3)   {$R_2$};
  \node[right,font=\footnotesize]       at (4.74,3.9)   {$R_3$};
  \node[right,font=\footnotesize]       at (4.74,2.3)   {$R_4$};
  \node[left,font=\footnotesize]        at (2.0,3.1)    {$B$};
  \node[right,font=\footnotesize]       at (4.6,3.1)    {$D$};

\end{circuitikz}}


\exname{直流双臂电桥} %实验名称
\instructor{宋斌} %指导教师
\class{电科2502} %班级
\name{管云哲} %姓名
\stuid{3250104519} %学号

\nyear{2026} %年
\nmonth{10} %月
\nday{13} %日
\nweekday{二} %星期几

\begin{document}
\setcounter{page}{1}
\makecover

\section{预习报告（10分）}
\subsection{实验综述（5分）}
（自述实验现象、实验原理和实验方法，包括必要的光路图、电路图、公式等。不超过500字。）

实验现象：电桥平衡时，检流计指针指零。

实验原理：用单臂电桥测低电阻时，接线电阻与接触电阻（$0.01\sim1\,\Omega$）会与待测电阻 $R_x$ 串联而被计入测量结果，误差极大。为此将低电阻改为四端接入：电流端 $C_{x1}$、$C_{x2}$ 通电流，电位端 $P_{x1}$、$P_{x2}$ 取电压，只有两电位端之间的电阻才是实测值，附加电阻被“转移”至待测电阻之外。把两个四端电阻 $R_x$、$R_s$ 接入桥路，并另设阻值均大于 $10\,\Omega$ 的标准电阻 $R_3$、$R_4$，用 $r<0.001\,\Omega$ 的粗导线连接 $R_x$ 与 $R_s$，即构成双臂电桥，其电路如图所示。

\begin{figure}[H]
    \centering
    \ArmBridge
    \caption{双臂电桥原理电路图}
    \label{fig:armbridge}
\end{figure}

平衡时 $I_G=0$，由基尔霍夫定律可得
\begin{equation}
    R_x=\frac{R_1}{R_2}R_s+\frac{R_4 r}{R_3+R_4+r}\left(\frac{R_1}{R_2}-\frac{R_3}{R_4}\right)
    \label{eq:rx}
\end{equation}
若使两对比率臂始终保持 $R_1/R_2=R_3/R_4$，则第二项为零，待测低电阻为
\begin{equation}
    R_x=\frac{R_1}{R_2}R_s
    \label{eq:rx2}
\end{equation}
金属的电阻随温度线性变化：$R_t=R_0(1+\alpha t)$，$R_0$ 为 $0\,^\circ\mathrm{C}$ 时的阻值，$\alpha$ 为电阻温度系数。

实验方法：（1）测电阻率。接通检流计并调零，灵敏度置于最低；接入待测金属导体，选适当倍率，先按 B 通电、后按 G 检流，调节粗调与细调盘使检流计偏转减小，逐步提高灵敏度至最高档调平衡，读数 $R_x=$ 倍率 $\times$（粗调$+$细调）；再用游标卡尺测导体直径 $d$、用米尺测长度 $L$。（2）测电阻温度系数。在最高灵敏度下调平衡，开启加热装置设定温度，用升温或降温法每改变约 $5\,^\circ\mathrm{C}$ 测一次 $R$ 与 $t$，记录 10 组数据。

\subsection{实验重点（3分）}
（简述本实验的学习重点，不超过100字。）

掌握双臂电桥测低电阻的原理，理解四端接法与双比率臂使 $R_1/R_2=R_3/R_4$ 以消除引线电阻、接触电阻影响的方法；学会 QJ44 型直流双臂电桥的操作及由 $R=\dfrac{\pi d^2}{4L}\rho$ 测定金属电阻率与电阻温度系数。

\subsection{实验难点（2分）}
（简述本实验的实现难点，不超过100字。）

操作顺序须“先按 B 后按 G”，并随平衡程度逐步提高灵敏度，否则检流计易过载；高灵敏度下平衡点难以判准，需兼顾调节速率与读数精度；升温测量时导体温度与显示温度不一致，须待温度稳定后同步读取温度与电阻。

{\fangsong 注意事项：
\begin{enumerate}[label=\arabic*.]
    \item 用PDF格式上传“实验报告”，文件名：学生姓名+学号+实验名称+周次。
    \item “实验报告”必须递交在“学在浙大”本课程内对应实验项目的“作业”模块内。
    \item “实验报告”成绩必须在“浙江大学物理实验教学中心网站”-“选课系统”内查询。
    \item 教学评价必须在“浙江大学物理实验教学中心网站”-“选课系统”内进行，学生必须进行教学评价，才能看到实验报告成绩，教学评价须在本次实验结束后3天内进行。
\end{enumerate}}
\begin{center}
    \fbfangsong{浙江大学物理实验教学中心制}
\end{center}
\end{document}
